\documentclass[11pt,a4paper]{article}

\usepackage[T1]{fontenc}
\usepackage[utf8]{inputenc}
\usepackage[spanish]{babel}
\usepackage{amsmath,amssymb,bm}
\usepackage{geometry}
\usepackage{booktabs}
\usepackage{array}
\usepackage{enumitem}
\usepackage{xcolor}
\usepackage{listings}
\usepackage{hyperref}
\usepackage{setspace}

\geometry{margin=2.2cm}
\setlist[enumerate]{leftmargin=*, itemsep=0.25em}
\setlist[itemize]{leftmargin=*, itemsep=0.2em}
\hypersetup{colorlinks=true, linkcolor=black, urlcolor=blue}

\lstset{
  basicstyle=\ttfamily\small,
  keywordstyle=\bfseries,
  commentstyle=\itshape\color{black!50},
  showstringspaces=false,
  breaklines=true,
  frame=single,
  rulecolor=\color{black!20},
  backgroundcolor=\color{black!3},
  columns=fullflexible,
  keepspaces=true,
}

\title{Actividad: estilo eruptivo con el conducto 1D estacionario}
\author{MA-5307 --- Análisis numérico de EDP \\
Departamento de Ingeniería Matemática, Universidad de Chile}
\date{}

\begin{document}
\maketitle
\thispagestyle{empty}

\noindent
Modelo de Kozono \& Koyaguchi (2009), marchado en diferencias finitas.
\textbf{No hace falta} \texttt{scikits.odes}. Notebook:
\texttt{actividad\_conducto1D.ipynb}.

\section{Qué hay que hacer}

\begin{enumerate}
  \item \textbf{No editen} \texttt{RIconduit1D\_FD.py},
        \texttt{RIconduitex5\_5.py} ni \texttt{RIconduitef5\_5.py}.
  \item Los datos de cada volcán están en \texttt{datos/}.
        Ahí sí pueden (y deben) escribir.
  \item Abran \texttt{actividad\_conducto1D.ipynb} y ejecuten todas las celdas.
        Comenten los gráficos en las celdas que dicen \emph{Comenten aquí}.
  \item El caso estromboliano lo \textbf{buscan ustedes}: copian la plantilla,
        rellenan y vuelven a correr.
\end{enumerate}

\section{Las dos ramas del tiro}

El mismo sistema de ecuaciones se cierra de dos maneras.

\begin{center}
\begin{tabular}{@{}>{\raggedright\arraybackslash}p{0.22\textwidth}
                >{\raggedright\arraybackslash}p{0.48\textwidth}
                >{\raggedright\arraybackslash}p{0.22\textwidth}@{}}
\toprule
\textbf{Rama} & \textbf{Qué pide en la boca} & \textbf{Típico de} \\
\midrule
explosiva
  & gas sónico, o $P\approx P_{\mathrm{atm}}$ \textbf{con} fragmentación
    ($\varphi\ge\varphi_{\mathrm{crit}}$)
  & Pliniana / sub-Pliniana \\[0.6em]
efusiva
  & $P\approx P_{\mathrm{atm}}$ \textbf{sin} fragmentar
    ($\varphi\le\varphi_{\mathrm{crit}}$)
  & lava / domo \\
\bottomrule
\end{tabular}
\end{center}

\vspace{0.4em}
Un mismo volcán puede tener solución explosiva, efusiva, \textbf{las dos},
o \textbf{ninguna}. Eso también se comenta.

\section{Y lo estromboliano, ¿cómo entra?}

Este modelo es \textbf{estacionario}: un perfil $P(z)$, $\varphi(z)$, $u(z)$
que no cambia en el tiempo. Una erupción estromboliana real es una sucesión
de explosiones discretas; \textbf{eso no lo resolvemos}.

Lo que sí pueden hacer es buscar datos petrológicos de un volcán
estromboliano (Stromboli, Villarrica, Etna, Pacaya, \ldots) y preguntarle
al modelo: \emph{¿con estos números el conducto sale fragmentado, como lava,
las dos cosas, o ninguna?} Eso es lo que hay que interpretar en los gráficos.

\section{Cómo usarlo}

En una terminal, desde la carpeta del proyecto:
\begin{lstlisting}[language=bash]
pip install numpy matplotlib
\end{lstlisting}

Después abran el notebook y ejecuten \textbf{todas} las celdas.
Para cambiar de volcán no toquen el solver: cambian el archivo en
\texttt{datos/} o copian la plantilla. La llamada es siempre:

\begin{lstlisting}[language=Python]
from datos import cargar
from RIconduit1D_FD import RIconduit1D_FD_f

volc = cargar("calbuco2015")   # archivo en datos/
out  = RIconduit1D_FD_f(volc, rama="explosiva", verbose=False)
\end{lstlisting}

\texttt{rama} es \texttt{"explosiva"} o \texttt{"efusiva"}.

\section{Casos}

\subsection{Erupción explosiva --- Calbuco 2015}

Datos \textbf{dados} en \texttt{datos/calbuco2015.py} (no los busquen).
Corran la rama explosiva.

\begin{center}
\small
\begin{tabular}{@{}llr@{}}
\toprule
& & Calbuco 2015 \\
\midrule
Estilo observado & & sub-Pliniana / explosiva \\
$\mathrm{SiO}_2$ (vidrio) & wt\% & $62.72$ \\
$\mathrm{H}_2\mathrm{O}$ & wt\% & $4$ \\
$T$ & $^\circ\mathrm{C}$ & $970$ \\
$\xi$ (cristales) & --- & $0.25$ \\
$H$ (cámara) & m & $-7000$ \\
Geometría & & cilindro \\
Radio $R$ & m & $16$ \\
$\Delta P$ & MPa & $5$ \\
\bottomrule
\end{tabular}
\end{center}

En los gráficos fíjense en:
\begin{itemize}
  \item dónde cruza $\varphi$ a $\varphi_{\mathrm{crit}}$
        (profundidad de fragmentación);
  \item si $u_g$ se separa de $u_m$ después de fragmentar;
  \item $P$ en la boca frente a $P_{\mathrm{atm}}$;
  \item el orden de magnitud del flujo de masa $Q$.
\end{itemize}

\paragraph{Comenten.}
¿A qué profundidad fragmenta? ¿Las velocidades se separan?
¿El flujo de masa es alto o bajo comparado con una lava?
¿La boca queda sónica o solo despresurizada?

\subsection{Erupción efusiva --- St.\ Helens 2004 (domo)}

Datos \textbf{dados} en \texttt{datos/sthelens2004.py}.
Corran la rama \textbf{efusiva}. Si quieren, corran también la explosiva
para ver si existe.

\begin{center}
\small
\begin{tabular}{@{}llr@{}}
\toprule
& & St.\ Helens 2004 \\
\midrule
Estilo observado & & domo / efusiva \\
$\mathrm{SiO}_2$ (vidrio) & wt\% & $73.5$ \\
$\mathrm{H}_2\mathrm{O}$ & wt\% & $3$ \\
$T$ & $^\circ\mathrm{C}$ & $950$ \\
$\xi$ (cristales) & --- & $0.40$ \\
$H$ (cámara) & m & $-5000$ \\
Geometría & & cilindro \\
Radio $R$ & m & $12.11$ \\
$\Delta P$ & MPa & $16.16$ \\
\bottomrule
\end{tabular}
\end{center}

Comparen con Calbuco: $\varphi$ vs $\varphi_{\mathrm{crit}}$,
$u_m$ y $u_g$, $P$ de salida, $Q$.

\paragraph{Comenten.}
¿Fragmenta? ¿$u_g$ se parece a $u_m$?
¿El tiro llega a $P_{\mathrm{atm}}$ o se queda corto?
Si la rama explosiva no llega a la boca, ¿qué significa eso para el estilo?

\subsection{Caso estromboliano --- lo buscan ustedes}

Este modelo \textbf{no} calcula las explosiones de Stromboli.
Calcula un estado estacionario. La pregunta es: con datos reales de un
volcán estromboliano, ¿qué rama(s) aparecen?

\paragraph{Pasos.}
\begin{enumerate}
  \item Copien \texttt{datos/plantilla\_estromboliana.py} a
        \texttt{datos/mi\_estromboliana.py}.
  \item Rellenen solo lo marcado con \texttt{\# BUSCAR}
        (composición del vidrio, $\mathrm{H}_2\mathrm{O}$, $T$, cristales,
        profundidad de cámara, radio o dique, sobrepresión).
  \item Si el magma es basáltico ($\mathrm{SiO}_2\lesssim 57\,\mathrm{wt\%}$),
        descomenten $C_1$ y $\beta$ basálticos en la plantilla
        ($C_1=6.80172\times 10^{-8}$, $\beta=0.7$).
        Si es silíceo, dejen Henry como está
        ($C_1=4.11\times 10^{-6}$, $\beta=0.5$).
  \item Pongan \texttt{LISTO = True} y citen las fuentes (paper, tabla, página).
  \item En el notebook, dejen
        \texttt{VOLCAN\_ESTROMBOLIANA = "mi\_estromboliana"}
        y vuelvan a ejecutar.
\end{enumerate}

Volcanes razonables: \textbf{Stromboli}, \textbf{Villarrica},
\textbf{Etna}, \textbf{Pacaya}.
Sobrepresión y radio casi nunca salen en un abstract: si no los encuentran,
elijan un valor en el rango de la plantilla y \textbf{justifiquen}.

Referencia útil:
Castruccio, Rebolledo \& Gómez (2025), \emph{JGR Solid Earth}
(Villarrica aparece como hawaiana/estromboliana).

Campos que tienen que buscar:

\begin{center}
\small
\begin{tabular}{@{}llp{0.48\textwidth}@{}}
\toprule
Variable & Unidades & Qué buscar \\
\midrule
composición del vidrio & wt\% & $\mathrm{SiO}_2$, $\mathrm{TiO}_2$, \ldots \\
\texttt{h2o1} & wt\% & $\mathrm{H}_2\mathrm{O}$ (inclusiones; típ.\ $0.5$--$2.5$) \\
\texttt{Tc1} & $^\circ\mathrm{C}$ & $T$ (basalto: $\sim 1050$--$1150$) \\
\texttt{xi1} & --- & fracción de cristales ($0$--$0.4$) \\
\texttt{H} & m & profundidad de cámara (negativa, ej.\ $-3000$) \\
\texttt{geometry} & --- & \texttt{"cylinder"} o \texttt{"dyke"} \\
\texttt{radius1} & m & radio, o semi-ancho si es dique \\
\texttt{overP1} & Pa & sobrepresión; si no aparece, prueben $1$--$5\,\mathrm{MPa}$ \\
\bottomrule
\end{tabular}
\end{center}

\paragraph{Comenten.}
¿Qué volcán eligieron y de dónde salieron los datos?
¿El modelo da rama explosiva, efusiva, ambas o ninguna?
¿Eso calza con el estilo observado, o el modelo estacionario se queda corto?
Si cambiaron radio o $\Delta P$ porque no los encontraron,
¿hacia qué lado mueve el resultado?

\section{Comparación}

Completen (una frase cada uno):
\begin{itemize}
  \item Calbuco vs St.\ Helens: la diferencia más clara en los perfiles es \ldots
  \item En este modelo, ``estromboliano'' se ve como \ldots
  \item Una limitación del 1D estacionario para el caso~3 es \ldots
\end{itemize}

\section*{Qué entregar}

El notebook con las figuras y los comentarios escritos,
más \texttt{datos/mi\_estromboliana.py} con las fuentes citadas.

\vspace{1.2em}
{\small
\noindent\textit{Datos de Calbuco y St.\ Helens:}
Castruccio, A., Rebolledo, A.\ \& Gómez, I.\ (2025).
The influence of melt composition, temperature, crystallinity and water
content on eruptive style and eruption rate: insights from a conduit model
of magma ascent. \emph{Journal of Geophysical Research: Solid Earth}, 130,
e2024JB030599. \url{https://doi.org/10.1029/2024JB030599}
}

\end{document}
